\documentclass[../main]{subfiles}
\begin{document}
\chapter{Principio 1}
\img{images/03/tractor.jpg}{Máquina de vapor agrícola.}
\info{
Contenido}{
  Primer principio de la termodinámica.
}
\actividad{ Investigar y contestar.}
\begin{enumerate}
  \item Enuncie el primer principio de la termodinámica.
  \item Si se aplica el primer principio a un sistema cerrado, ¿Que se obtiene?
  \item Definir el equivalente mecánico del calor, dar fórmula.
  \item Defina entropía.
  \item Defina entalpía.
  \item Explique el primer principio aplicado a las transformaciones abiertas.
\end{enumerate}

\actividad{Resolver los siguientes ejercicios.}
\begin{enumerate}
  \item Calcular la cantidad de calor que es necesario entregar a un
    gas para que realice un trabajo de expansión igual a $W=2562kgm$
    suponiendo que su energía interna es constante.
  \item Calcular el trabajo externo que es necesario ejercer sobre un
    gas cuando al comprimirlo entrega $30kcal$ de energía interna constante.
  \item Calcular la variación de energía interna que experimenta un
    gas cuando es comprimido a la presión de $p=2atm$ mediante el
    desplazamiento de un émbolo de $A=0,01m^2$ de area transversal y
    $d=0,3m$ de carrera suponiendo que durante la compresión pasan
    $5kcal$ de calor desde el cilindro al recipiente de refrigeración.
  \item Calcular la variación de energía interna que experimenta una
    masa gaseosa cuando su volumen aumenta $\Delta V=0,5m^3$,
    manteniéndose una presión constante de $P=4atm$, en el supuesto
    de que durante la transformación no hay intercambio de calor.
  \item Calcular la cantidad de calor que absorbe un gas que recorre
    la transformación $1\cdot2\cdot3$ de la figura realizando un
    trabajo externo de $1000kgm$, si al recorrer la transformación
    $1\cdot3$ absorbe $800kcal$ y realiza un trabajo de $50000kgm$

    \begin{center}
      \begin{tikzpicture}
        \draw[-latex] (-0.5,0) -- (4,0) node[below] {$V$};
        \draw[-latex] (0,-0.5) -- (0,4) node[left] {$P$};
        \draw[latex-] (1,1) node [left]{1} -- (3,1);
        \draw[latex-] (3,1) node [right]{3} -- (3,3);
        \draw[latex-] (3,3) node [right]{2} -- (1,1);
      \end{tikzpicture}
    \end{center}

  \item Un compresor aspira aire a la presión $p_i=1atm$ y con  peso
    específico $\gamma_{i}=1,25kg/m^3$ y lo expulsa a la presión de
    $p_e=5atm$, con el peso específico $\gamma_e=4kg/m^3$. Suponiendo
    que la energía interna pasó de $4kcal/kg$ a $32 Kcal/kg$.
    Calcular la diferencia de entalpía por unidad de peso que
    experimentó la masa de aire.

\end{enumerate}
\end{document}

